%%%%%%%%%%  scan idx 444  |  folio 431  %%%%%%%%%%
alors il en est de même pour \(\bigoplus_{i\ge0}\Tor_i^A(B,C)\).

Cette dernière remarque s'applique au cas où \(B\) est le quotient de \(A\) par un idéal régulier, car on peut alors prendre pour \(P_\bullet\) le complexe de Koszul associé à un homomorphisme régulier définissant l'idéal de \(B\).

\subsection*{2.3}
La structure de \(B\otimes_A C\)-algèbre définie en 2.1 sur \(\bigoplus_{i\ge0}\Tor_i^A(B,C)\) est fonctorielle en \(B\) et \(C\). On peut alors globaliser les constructions précédentes: soient \(X\) un schéma, \(X'\), \(X''\) deux \(X\)-schémas; alors on peut localement sur \(X'\times_X X''\) définir les morphismes (2.1.1), ils se recollent par fonctorialité, et définissent sur \(\bigoplus_{i\ge0}\Tor_i^{\cO_X}(\cO_{X'},\cO_{X''})\) une structure de \(\cO_{X'\times_X X''}\)-Algèbre graduée associative.

\textbf{Lemme 2.4.} Soient \(X',X''\) deux \(X\)-schémas, et \(M\) un \(\cO_{X'}\)-Module plat. Alors pour tout \(i\ge0\), on a un isomorphisme canonique :
\begin{equation}\tag{2.4.1}
\Tor_i^{\cO_X}(M,\cO_{X''})\simeq
M\otimes_{\cO_{X'}}\Tor_i^{\cO_X}(\cO_{X'},\cO_{X''}).
\end{equation}
Il suffit de définir un isomorphisme fonctoriel dans le cas affine; on pourra alors par recollement le définir globalement dans le cas général. Soient alors \(X=\Spec(A)\), \(X'=\Spec(A')\), \(X''=\Spec(A'')\); \(M\) est défini par un \(A'\)-module plat \(M\). Soit \(P_\bullet\) une résolution projective de \(A''\) sur \(A\); les \(\Tor_i^A(A',A'')\) sont les groupes d'homologie de \(A'\otimes_A P_\bullet\). Comme \(M\) est plat sur \(A'\), on a un isomorphisme canonique : \(M\otimes_{A'}\Tor_i^A(A',A'')\xrightarrow{\sim}H_i\bigl(M\otimes_{A'}(A'\otimes_A P_\bullet)\bigr)\). Par ailleurs, l'isomorphisme canonique : \(M\otimes_{A'}(A'\otimes_A P_\bullet)\xrightarrow{\sim}M\otimes_A P_\bullet\) donne un isomorphisme canonique : \(H_i\bigl(M\otimes_{A'}(A'\otimes_A P_\bullet)\bigr)\xrightarrow{\sim}\Tor_i^A(M,A'')\), d'où le résultat.

\textbf{Proposition 2.5.} Soient \(X\) un schéma, \(Y\) et \(Z\) deux sous-schémas fermés de \(X\), \(T=Y\times_X Z\). On suppose que pour tout \(x\in T\), il existe un voisinage ouvert \(U\) de \(x\), et une
%%%%%%%%%%  scan idx 445  |  folio 432  %%%%%%%%%%
suite de sections \(f_1,\ldots,f_p,g_1,\ldots,g_r,h_1,\ldots,h_s\) de \(\cO_U\) telles que le complexe de Koszul associé à \(f_1,\ldots,f_p,g_1,\ldots,g_r\) soit une résolution de \(\cO_{Y\cap U}\), le complexe de Koszul associé à \(f_1,\ldots,f_p,h_1,\ldots,h_s\) une résolution de \(\cO_{Z\cap U}\), et le complexe de Koszul associé à \(f_1,\ldots,f_p,g_1,\ldots,g_r,h_1,\ldots,h_s\) une résolution de \(\cO_{T\cap U}\) (conditions satisfaites en particulier dans le cas noethérien si la suite \(f_1,\ldots,f_p,g_1,\ldots,g_r,h_1,\ldots,h_s\) est régulière et si \(Y=V(f_1,\ldots,f_p,g_1,\ldots,g_r)\), \(Z=V(f_1,\ldots,f_p,h_1,\ldots,h_s)\)). Soient \(I,J\) les idéaux de \(Y,Z\). On a alors les isomorphismes suivants:
\[
\begin{array}{ll}
\text{i)}& \forall i\ge0,\quad
\Tor_i^{\cO_X}(\cO_Y,\cO_Z)\xrightarrow{\sim}
\bigwedge^i\Tor_1^{\cO_X}(\cO_Y,\cO_Z);\\[0.8ex]
\text{ii)}&
\Tor_1^{\cO_X}(\cO_Y,\cO_Z)\simeq
I\cap J/I.J\simeq
\Ker\bigl(i^*(N_{Y/X})\times j^*(N_{Z/X})\to N_{T/X}\bigr),
\end{array}
\]
où on note \(i\) l'immersion de \(T\) dans \(Y\) et \(j\) l'immersion de \(T\) dans \(Z\).

En particulier, pour toute immersion fermée régulière \(Y\to X\), on a
\[
\begin{array}{ll}
\text{i')}& \forall i\ge0,\quad
\Tor_i^{\cO_X}(\cO_Y,\cO_Y)\xrightarrow{\sim}
\bigwedge^i\Tor_1^{\cO_X}(\cO_Y,\cO_Y);\\[0.8ex]
\text{ii')}&
\Tor_1^{\cO_X}(\cO_Y,\cO_Y)\simeq N_{Y/X}.
\end{array}
\]

\emph{Démonstration.} i) On a vu en 2.3 que \(\bigoplus_{i\ge0}\Tor_i^{\cO_X}(\cO_Y,\cO_Z)\) pouvait être muni d'une structure de \(\cO_T\)-Algèbre; d'après 2.2, cette algèbre est alternée, car l'immersion de \(Y\) (ou \(Z\)) dans \(X\) est régulière. Il existe donc un homomorphisme d'algèbres et un seul:
\begin{equation}\tag{2.5.1}
\bigwedge \Tor_1^{\cO_X}(\cO_Y,\cO_Z)\longrightarrow
\bigoplus_{i\ge0}\Tor_i^{\cO_X}(\cO_Y,\cO_Z),
\end{equation}
prolongeant l'identité en degré \(1\).

Pour voir que c'est un isomorphisme, on peut le voir localement. On se place donc au voisinage d'un point de \(T\), et on considère des sections \(f_1,\ldots,f_p,g_1,\ldots,g_r,h_1,\ldots,h_s\) de \(\cO_X\) dans ce voisinage, ayant les propriétés indiquées dans l'énoncé de 2.5. Posons alors \(E=(\cO_X)^p\), \(E'=(\cO_X)^r\), et soient \(u:E\to\cO_X\) l'homomorphisme défini par les \(f_i\), \(v:E'\to\cO_X\) l'homomorphisme défini par les \(g_j\), \(K_\bullet(u)\) et \(K_\bullet(v)\) les
%%%%%%%%%%  scan idx 446  |  folio 433  %%%%%%%%%%
complexes de Koszul associés. Le complexe de Koszul associé à la suite \(f_1,\ldots,f_p,g_1,\ldots,g_r\) est isomorphe au complexe produit tensoriel \(K_\bullet(u)\otimes_{\cO_X}K_\bullet(v)\). Comme c'est une résolution de \(\cO_Y\) par des \(\cO_X\)-Modules libres de type fini, les \(\Tor_i^{\cO_X}(\cO_Y,\cO_Z)\) sont les faisceaux d'homologie du complexe \(K_\bullet(u)\otimes_{\cO_X}K_\bullet(v)\otimes_{\cO_X}\cO_Z\). On peut le considérer comme le produit tensoriel des complexes \(K_\bullet(u)\) et \(K_\bullet(v)\otimes_{\cO_X}\cO_Z\), et utiliser pour calculer son homologie les suites spectrales associées à un bicomplexe (EGA \(0_{\mathrm{III}}\) (11.3.5)), qui ont comme aboutissement l'homologie du complexe simple associé.

On considère la suite spectrale telle que \({}'E^2_{pq}=H^I_p\bigl(H^{II}_q(K_{\bullet,\bullet})\bigr)\), avec les notations de EGA \(0_{\mathrm{III}}\) (11.3.5), le premier indice étant relatif au complexe \(K_\bullet(u)\) et le second au complexe \(K_\bullet(v)\otimes_{\cO_X}\cO_Z\). Or le complexe \(K_\bullet(v)\otimes_{\cO_X}\cO_Z\) est le complexe de Koszul sur \(\cO_Z\) associé à la suite de sections \(\overline g_1,\ldots,\overline g_r\), où les \(\overline g_i\) sont les images des \(g_i\) dans \(\cO_Z\), et ce complexe est une résolution de \(\cO_T\) sur \(\cO_Z\) car les immersions \(Z\to X\) et \(T\to Z\to X\) sont régulières (cf. démonstration de 1.7 ii). Comme \(\bigwedge^p E\) est libre sur \(\cO_X\), le complexe \(\bigwedge^pE\otimes_{\cO_X}(K_\bullet(v)\otimes_{\cO_X}\cO_Z)\) est une résolution de \(\bigwedge^pE\otimes_{\cO_X}\cO_T\). Par conséquent la suite spectrale dégénère, et les \(\Tor_i^{\cO_X}(\cO_Y,\cO_Z)\) sont les faisceaux d'homologie du complexe \(K_\bullet(u)\otimes_{\cO_X}\cO_T\). Or la différentielle de \(K_\bullet(u)\) s'exprime en termes des \(f_i\), qui annulent \(\cO_T\); par suite, la différentielle de \(K_\bullet(u)\otimes_{\cO_X}\cO_T\) est nulle, et on obtient donc les isomorphismes:
\begin{equation}\tag{2.5.2}
\Tor_i^{\cO_X}(\cO_Y,\cO_Z)
\xrightarrow{\sim}
\cO_T\otimes_{\cO_X}\bigwedge^iE
\xrightarrow{\sim}
\bigwedge^i(\cO_T)^p.
\end{equation}
En particulier, on trouve pour \(i=1\):
\[
\Tor_1^{\cO_X}(\cO_Y,\cO_Z)\xrightarrow{\sim}(\cO_T)^p,
\]
d'où pour tout \(i\):
\begin{equation}\tag{2.5.3}
\Tor_i^{\cO_X}(\cO_Y,\cO_Z)
\xrightarrow{\sim}
\bigwedge^i\Tor_1^{\cO_X}(\cO_Y,\cO_Z).
\end{equation}

%%%%%%%%%%  scan idx 447  |  folio 434  %%%%%%%%%%
Or il
résulte de 2.2 que l'isomorphisme
\[
\bigwedge(\cO_T)^n\xrightarrow{\sim}
\bigoplus_{i\ge0}\Tor_i^{\cO_X}(\cO_Y,\cO_Z)
\]
défini par les isomorphismes (2.5.2) est compatible aux structures d'algèbres des deux membres; il en est donc de même pour l'isomorphisme
\[
\bigwedge\Tor_1^{\cO_X}(\cO_Y,\cO_Z)\xrightarrow{\sim}
\bigoplus_{i\ge0}\Tor_i^{\cO_X}(\cO_Y,\cO_Z)
\]
défini par les isomorphismes (2.5.3). Pour des raisons d'unicité, ce dernier est donc la restriction de l'homomorphisme (2.5.1) qui avait été défini globalement. L'homomorphisme (2.5.1) est donc un isomorphisme.

ii) On considère la suite exacte:
\[
0\longrightarrow I\longrightarrow \cO_X\longrightarrow \cO_Y\longrightarrow0,
\]
et la suite exacte des Tor qui en résulte:
\[
0\longrightarrow\Tor_1^{\cO_X}(\cO_Y,\cO_Z)
\longrightarrow I\otimes_{\cO_X}\cO_Z
\longrightarrow \cO_Z.
\]
Il en résulte que \(\Tor_1^{\cO_X}(\cO_Y,\cO_Z)\) est le noyau de l'homomorphisme :
\[
I\otimes_{\cO_X}(\cO_X/J)\longrightarrow \cO_X/J,
\]
c'est-à-dire \(I\cap J/I.J\).

Pour montrer la dernière relation, nous allons montrer qu'il existe une suite exacte:
\[
0\longrightarrow I\cap J/I.J
\longrightarrow i^*(N_{Y/X})\times j^*(N_{Z/X})
\longrightarrow N_{T/X}\longrightarrow0.
\]
Les homomorphismes \(i^*(N_{Y/X})\to N_{T/X}\) et \(j^*(N_{Z/X})\to N_{T/X}\) ont été définis en (1.6.1); on prend ici leur différence. Par ailleurs, \(I\cap J\) s'envoie de façon naturelle dans \(I/I^2\), donc dans \(i^*(I/I^2)\), et de même dans \(j^*(J/J^2)\), et ces homomorphismes s'annulent sur \(I.J\), car \(I,J\subset K\) (où \(K\) désigne l'idéal de \(T\)) ce qui définit la flèche de gauche.

%%%%%%%%%%  scan idx 448  |  folio 435  %%%%%%%%%%
Pour vérifier que la suite est exacte, on peut le faire localement sur \(T\), et utiliser les sections définies dans l'énoncé de 2.5. On a alors une base de \(N_{T/X}\) formée des images des \(f_i,g_j,h_k\), une base de \(i^*(N_{Y/X})\) formée des images des \(f_i,g_j\), et une base de \(j^*(N_{Z/X})\) formée des images des \(f_i,h_k\). Il est alors clair que \(i^*(N_{Y/X})\times j^*(N_{Z/X})\to N_{T/X}\) est surjectif, et que son noyau est formé des éléments ayant des coordonnées nulles sur les \((\overline g_j,0)\) et les \((0,\overline h_k)\), et des coordonnées égales pour tout \(i\) sur \((\overline f_i,0)\) et \((0,\overline f_i)\). Par suite, la suite est également exacte au milieu. Or les trois \(\cO_T\)-Modules intervenant sont localement libres de type fini; comme \(\Tor_1^{\cO_X}(\cO_Y,\cO_Z)\) est de rang \(p\) d'après la démonstration du i), on voit immédiatement que la suite est exacte.

\subsection*{2.6}
On suppose maintenant que \(X\) est un schéma ayant un faisceau ample. Tout sous-schéma fermé \(Y\) de \(X\) a alors un faisceau ample. On sait que dans ce cas ``il existe des résolutions globales'', et que les \(K^\circ\) formés respectivement avec les faisceaux localement libres de type fini et les complexes parfaits sont identiques (II, 2.2.9). Si en outre l'immersion \(i:Y\to X\) est régulière, c'est un morphisme parfait, et on peut définir un homomorphisme additif \(i_*:K^\circ(Y)\longrightarrow K^\circ(X)\) (IV 2.12), en associant à la classe d'un complexe parfait \(L^\bullet\) la classe du complexe \(Ri_*(L^\bullet)=i_*(L^\bullet)\) (\(i_*\) étant un foncteur exact).

\textbf{Proposition 2.7.} Soient \(X\) un schéma ayant un faisceau ample, \(i:Y\to X\) une immersion fermée régulière, \(N\) le faisceau conormal de \(Y\) dans \(X\). Alors, pour tout \(x\in K^\circ(Y)\), on a la relation:
\begin{equation}\tag{2.7.1}
i^*(i_*(x))=x.\lambda_{-1}(N),
\end{equation}
où \(N\) est la classe de \(N\) dans \(K^\circ(Y)\), et \(\lambda_{-1}(N)=\sum_{i\ge0}(-1)^i\lambda^i(N)\).

\emph{Démonstration.}
%%%%%%%%%%  scan idx 449  |  folio 436  %%%%%%%%%%
